\documentclass[12pt]{article}

\usepackage{sbc-template}
\usepackage{graphicx,url}
\usepackage[utf8]{inputenc}
\usepackage[brazil]{babel}
\usepackage{multirow}
\usepackage{amsmath}

\sloppy

\title{Análise Comparativa de Protocolos de Roteamento AODV vs OLSR em Redes MESH Escaláveis: Investigação do Ponto de Transição para Cenários de Smart Cities}

\author{Aluno UFABC\inst{1}, Prof. Orientador\inst{1}}

\address{Centro de Engenharia, Modelagem e Ciências Sociais Aplicadas (CECS)\\
  Universidade Federal do ABC (UFABC) -- Santo André -- SP -- Brazil
  \email{aluno@ufabc.edu.br, orientador@ufabc.edu.br}
}

\begin{document}

\maketitle

\begin{abstract}
  Smart cities and Internet of Things applications increasingly rely on wireless mesh networks for distributed communication. The choice between reactive (e.g., AODV) and proactive (e.g., OLSR) routing protocols significantly impacts network performance at different scales. This paper presents a comprehensive comparative analysis of AODV and OLSR in mesh networks ranging from 100 to 1500 nodes through simulation-based experiments using MeshSim (NS-3). We evaluate five key metrics: packet delivery ratio (PDR), end-to-end latency, routing overhead, convergence time, and CPU utilization. Our results reveal the scalability characteristics and identify critical transition points where one protocol becomes more suitable than the other. The findings provide practical recommendations for network designers deploying mesh networks in smart city applications.
\end{abstract}

\begin{resumo}
  Cidades inteligentes e aplicações de Internet das Coisas dependem cada vez mais de redes mesh sem fio para comunicação distribuída. A escolha entre protocolos de roteamento reativos (ex: AODV) e proativos (ex: OLSR) impacta significativamente o desempenho em diferentes escalas de rede. Este artigo apresenta uma análise comparativa abrangente de AODV e OLSR em redes mesh variando de 100 até 1500 nós através de experimentos baseados em simulação usando MeshSim (NS-3). Avaliamos cinco métricas principais: taxa de entrega de pacotes (PDR), latência fim-a-fim, overhead de roteamento, tempo de convergência e utilização de CPU. Nossos resultados revelam as características de escalabilidade e identificam pontos críticos de transição onde um protocolo é mais adequado que o outro. Os achados fornecem recomendações práticas para engenheiros implantando redes mesh em aplicações de cidades inteligentes.
\end{resumo}

\section{Introdução}

Redes mesh sem fio multi-hop (Multi-hop Wireless Mesh Networks) tornaram-se fundamentais em cenários de IoT e cidades inteligentes, onde dispositivos comunicam-se de forma autônoma sem infra-estrutura centralizada. Essas redes apresentam relevância especial em ambientes em que a instalação de cabeamento tradicional é impraticável, como em monitoramento ambiental distribuído, iluminação inteligente urbana e sistemas de rastreamento de ativos \cite{wireless_akyildiz_2005,wireless_sichitiu_2004}.

O roteamento em redes mesh é crítico pois determina latência, consumo energético, confiabilidade e escalabilidade da rede. Os protocolos de roteamento em WMNs classificam-se em três categorias: proativos, reativos e híbridos \cite{wireless_benyamina_2012}. Os protocolos proativos, também denominados orientados a tabela (\textit{table-driven}), estabelecem caminhos para todos os nós destino acessíveis de forma contínua e independente da demanda imediata, garantindo informações de roteamento sempre atualizadas e baixa latência de entrega da primeira mensagem. Os protocolos reativos, ou sob demanda (\textit{on-demand}), estabelecem rotas apenas quando necessário: o processo de descoberta de rota é iniciado pelo nó origem e prossegue até que um caminho seja encontrado ou todas as alternativas se esgotem, reduzindo o overhead de controle em redes com baixa mobilidade. Cada abordagem apresenta trade-offs distintos em termos de overhead de controle, latência de descoberta de rotas e adaptabilidade a mudanças topológicas.

Contudo, falta investigação empírica rigorosa em escala contemporânea sobre como estes protocolos se comportam quando redes crescem de centenas para milhares de nós. Recomendações atuais em literatura baseiam-se frequentemente em benchmarks desatualizados ou estudos teóricos sem validação experimental robusta. Este cenário deixa engenheiros de smart cities sem dados confiáveis para tomar decisões arquiteturais.

O presente trabalho endereça esta lacuna por meio de uma análise comparativa controlada de AODV versus OLSR em escalas que refletem aplicações reais de smart cities. Utiliza-se simulação em MeshSim (baseada em NS-3) para garantir reprodutibilidade absoluta e coleta precisa de métricas. A contribuição principal contempla um benchmark quantitativo que identifica o ponto de transição crítico onde o overhead quadrático de OLSR começa a prejudicar significativamente o desempenho, enquanto AODV mantém escalabilidade linear.

\section{Trabalhos Relacionados}

O estudo de redes mesh sem fio e protocolos de roteamento é uma área madura em pesquisa. \cite{wireless_akyildiz_2005} e \cite{wireless_sichitiu_2004} estabeleceram os fundamentos conceituais de redes mesh, discutindo arquitetura, caracterização de falhas, e desafios de escalabilidade. Posteriormente, \cite{wireless_benyamina_2012} apresentou um survey abrangente sobre design de redes mesh, consolidando conhecimento sobre diferentes paradigmas de roteamento.

No contexto de avaliação de desempenho específica, \cite{mesh_passos_2005} conduziram medições práticas de performance em redes mesh, fornecendo baseline histórico importante. Trabalhos mais recentes focam em aplicações específicas: \cite{research_lee_2024} investigaram sistemas de comunicação mesh para detecção de fogo e gás perigoso, enquanto \cite{performance_khalifeh_2018} compararam redes DigiMesh e ZigBee visando avaliar trade-offs entre diferentes implementações.

Aplicações emergentes em cidades inteligentes motivam novas investigações. \cite{optimizing_ruete_2024} exploraram otimizações de redes mesh baseadas em ZigBee e DigiMesh para rotas de evacuação iluminadas em cenários de simulação de desastres, demonstrando a relevância contemporânea de redes mesh.

Nossas análises complementam este corpo de trabalho por meio de quatro contribuições principais. Primeiro, realiza-se uma avaliação controlada em escala de cem a mil e quinhentos nós com rigor estatístico garantido por simulação determinística. Segundo, apresenta-se comparação direta entre AODV e OLSR sob condições padronizadas, evitando os confundidores típicos de estudos anteriores. Terceiro, identificam-se pontos críticos de transição entre protocolos, permitindo que engenheiros tomem decisões informadas sobre qual protocolo adotar em cada regime de escala. Finalmente, oferece-se matriz de recomendação prática para engenheiros que implantam redes mesh em cenários reais de cidades inteligentes.

\section{Metodologia}

\subsection{Ambiente de Simulação}

Utilizamos MeshSim \cite{lorinder_github}, um simulador de redes mesh de código
aberto construído sobre o framework NS-3, adaptado e estendido para os
experimentos deste trabalho.
O código e os scripts de configuração e análise encontram-se disponíveis em
repositório público \cite{caburlao_github}.
O ambiente de simulação executa em Linux com bibliotecas NS-3 compiladas em modo
\textit{Release} (CMake, versão de desenvolvimento de abril de 2026),
garantindo desempenho determinístico e reprodutível.

A escolha de MeshSim sobre alternativas como OPNET ou OMNeT++ justifica-se por cinco fatores interrelacionados. O framework NS-3, que constitui o núcleo do MeshSim, oferece maturidade comprovada e validação empírica extensiva na comunidade de pesquisa. O suporte nativo e independente aos protocolos AODV e OLSR, sem necessidade de modificações de código, elimina vieses de implementação. Os modelos de canal Wi-Fi incorporam mecanismos calibráveis de atenuação por distância (Yans e LogDistance), permitindo modelagem realista do ambiente de propagação. A reprodução determinística é garantida por sementes aleatórias explícitas, assegurando que qualquer pesquisador reproduza exatamente os mesmos resultados. Finalmente, a integração nativa com o FlowMonitor possibilita coleta automatizada de métricas granulares por fluxo de dados.

\subsection{Topologia e Escalonamento de Rede}

Adota-se topologia aleatória com distribuição uniforme de nós, exceto pela topologia base de 25 nós que segue grade regular $5 \times 5$ com espaçamento de 50\,m entre nós adjacentes. Para todas as topologias aleatórias, a densidade de nós é mantida constante em $\rho = 2{,}84 \times 10^{-4}$\,nós/m$^2$, derivada da topologia de referência (50 nós em $420 \times 420$\,m).

A área de cada topologia é calculada como:
\begin{equation}
  A = \sqrt{\frac{N}{\rho}} \times \sqrt{\frac{N}{\rho}} \quad \text{(metros)}
\end{equation}

Investigam-se dez pontos de escala cobrindo a faixa de 25 a 1000 nós, conforme Tabela~\ref{tab:topologias}. Todas as topologias aleatórias utilizam distância mínima entre nós de 28\,m e alcance de transmissão de 115\,m (IEEE 802.11g, modelo LogDistance). A densidade constante garante que variações de desempenho reflitam o efeito da escala (diâmetro da rede, número de saltos) e não variações de cobertura. A arquitetura do MeshSim organiza-se em cinco camadas funcionais de nós, cada uma desempenhando papéis específicos na simulação. A primeira camada compreende os nós mesh, que formam o \textit{backbone} sem fio sob padrão IEEE 802.11g, utilizando canal Yans, controle de taxa ARF e protocolo Dot11s. 

\begin{table}[h]
\centering
\caption{Parâmetros das topologias simuladas.}
\label{tab:topologias}
\begin{tabular}{lrrrr}
\hline
\textbf{Topologia} & \textbf{Nós} & \textbf{Área (m)} & \textbf{Aquecimento (s)} & \textbf{Duração (s)} \\
\hline
grid-25      &   25 & $200 \times 200$   & 120 & 240 \\
random-50    &   50 & $420 \times 420$   & 180 & 300 \\
random-75    &   75 & $514 \times 514$   & 210 & 330 \\
random-100   &  100 & $600 \times 600$   & 240 & 360 \\
random-150   &  150 & $727 \times 727$   & 270 & 390 \\
random-200   &  200 & $840 \times 840$   & 300 & 420 \\
random-300   &  300 & $1028 \times 1028$ & 330 & 450 \\
random-500   &  500 & $1330 \times 1330$ & 360 & 480 \\
random-750   &  750 & $1626 \times 1626$ & 390 & 510 \\
random-1000  & 1000 & $1880 \times 1880$ & 420 & 540 \\
\hline
\end{tabular}
\end{table}

\subsection{Configuração dos Protocolos de Roteamento}

Ambos os protocolos encontram-se instalados com parâmetros padrão do NS-3.

O AODV (\textit{Ad hoc On-Demand Distance Vector}) é um protocolo reativo que facilita a criação de rotas multi-hop dinâmicas entre nós, estabelecendo-as apenas quando demandadas por solicitações de comunicação específicas. Mensagens de \textit{Route Request} (RREQ) são difundidas pelo nó origem; o destino ou um nó intermediário com rota válida responde com uma mensagem \textit{Route Reply} (RREP). Um aspecto central do AODV é o emprego de números de sequência de destino em cada entrada de rota: o próprio destino gera esse número e o compartilha com os nós solicitantes junto aos detalhes da rota. Esse mecanismo garante roteamento livre de laços (\textit{loop-free}) e elimina o problema de ``contagem ao infinito'' associado ao algoritmo de Bellman-Ford. Em caso de falha de enlace, o protocolo notifica prontamente os nós afetados para que invalidem as rotas comprometidas. O AODV é dimensionado para redes de dezenas a milhares de nós, priorizando a minimização do tráfego de controle e do overhead sobre os dados. Sua configuração no presente trabalho utiliza \texttt{HelloInterval=1s}, \texttt{ActiveRouteTimeout=3s}, \texttt{PathDiscoveryTime=3.5s} e \texttt{RreqRetries=2}.

O OLSR (\textit{Optimized Link State Routing Protocol}) implementa o paradigma proativo orientado a tabela (\textit{table-driven}), projetado especificamente para redes ad hoc móveis. O mecanismo central do OLSR é a seleção de \textit{Multipoint Relays} (MPRs): cada nó escolhe um subconjunto de seus vizinhos a um salto com conexões bidirecionais simétricas para atuar como MPR. Somente os MPRs retransmitem mensagens de controle de topologia (\textit{Topology Control} -- TC), reduzindo significativamente o número de retransmissões necessárias em comparação ao \textit{flooding} completo. O OLSR garante rotas de caminho mínimo para todos os destinos exigindo que os MPRs declarem informações de estado de enlace de seus próprios MPRs selecionados; informações de enlace adicionais podem ser utilizadas para redundância. As mensagens de controle são trocadas periodicamente, mantendo as tabelas de roteamento de todos os nós continuamente atualizadas. Sua configuração no presente trabalho utiliza \texttt{HelloInterval=2s}, \texttt{TcInterval=5s} e \texttt{MidInterval=5s}.

O tempo de aquecimento (\textit{warm-up}) de cada topologia é dimensionado para garantir convergência completa antes do início das medições, segundo critérios distintos para cada protocolo.

Para o OLSR, aplica-se o critério $T_{warmup} \geq D_{hops} \times T_{TC}$, onde $D_{hops}$ é o diâmetro estimado da rede e $T_{TC} = 5$\,s o intervalo de \textit{Topology Control}. Para redes de 1000 nós ($D_{hops} \approx 12$ saltos), o aquecimento de 420\,s
garante pelo menos 84 rodadas de TC propagadas pelo diâmetro completo, assegurando que todas as tabelas de roteamento estarão convergidas.

Para o AODV, o período de aquecimento é utilizado em conjunto com os \textit{probe flows} descritos na Seção~\ref{subsec:traffic}: fluxos de eco curtos enviados nos 30\,s anteriores ao início das medições forçam a descoberta de rota para todos os destinos. Como o \texttt{ActiveRouteTimeout} é de 3\,s e os probes terminam 1\,s antes das medições, as rotas permanecem válidas no cache AODV quando os fluxos de medição iniciam, garantindo que as métricas reflitam comportamento em regime permanente e não latência de primeira descoberta.

\subsection{Padrão de Tráfego}
\label{subsec:traffic}

Antes do início das medições, cada experimento executa uma fase de \textit{probe} com duração de 30\,s ($T_{warmup} - 30$\,s a $T_{warmup} - 1$\,s). Nessa fase, cinco fluxos UDP de eco de baixa taxa são enviados de cada wiredSta ao servidor backhaul, forçando o AODV a completar a descoberta de rotas para todos os destinos antes que os fluxos de medição comecem. Os fluxos de probe são excluídos da análise por dois critérios combinados: (i) filtro de instante de início --- apenas fluxos com $\texttt{timeFirstTxPacket} \geq T_{warmup}$ são considerados; e (ii) filtro direcional via \texttt{Ipv4FlowClassifier} --- somente fluxos no sentido wiredSta$\to$backhaul UDP são contabilizados, excluindo os \textit{echo replies} (backhaul$\to$wiredSta) e o fluxo TCP de transferência de arquivo. Esses critérios garantem que as métricas reportadas reflitam o comportamento em regime permanente (\textit{steady-state}) de ambos os protocolos sobre exatamente o mesmo conjunto de fluxos.

Cada experimento instância cinco fluxos UDP de eco independentes de medição, um por wiredSta. Os fluxos são iniciados de forma escalonada, com intervalo de 5\,s entre cada início ($T_{start,i} = T_{warmup} + (i-1) \times 5$\,s), evitando o caso patológico de $N$ descobertas de rota simultâneas. Adicionalmente, um fluxo TCP de transferência de arquivo é iniciado após o último fluxo echo, do servidor backhaul ($10.1.1.1$) para a wiredSta mais distante.

Cada cliente UDP envia 100 pacotes de 128 bytes a 2\,pacotes/s (\texttt{Interval=0.5s}), cobrindo uma janela de medição de 50\,s por fluxo. A duração total de cada simulação é fixada em $T_{warmup} + 105$\,s, contemplando 25\,s de escalonamento, 50\,s de medição e 30\,s de buffer para conclusão do fluxo TCP. O servidor de eco responde ao remetente com o mesmo payload, de forma que o FlowMonitor captura métricas nos dois sentidos (upload e echo de retorno).

\subsection{Métricas Coletadas}

As métricas são extraídas exclusivamente do FlowMonitor, em granularidade de fluxo. Apenas os fluxos UDP no sentido wiredSta para backhaul são incluídos no cálculo, conforme critérios descritos na Seção~\ref{subsec:traffic}.

\paragraph{Taxa de entrega de pacotes (PDR).}
O FlowMonitor registra \texttt{txPackets} somente para fluxos efetivamente enfileirados na camada IP. Protocolos proativos como OLSR, ao não possuírem rota para um destino inalcançável, impedem o enfileiramento e resultam em \texttt{txPackets}$=0$ para esse fluxo --- que é descartado pelo FlowMonitor.Protocolos reativos como AODV, por sua vez, enfileiram o pacote enquanto aguardam a conclusão do RREQ, contabilizando \texttt{txPackets} mesmo para destinos que eventualmente se revelam inalcançáveis.

Desta forma, para calcular a taxa de entrega de pacotes (PDR), utilizou-se $\sum tx_i$ como denominador tornaria a comparação assimétrica. Adota-se, portanto, um denominador fixo igual ao total de pacotes esperados:
\begin{equation}
  \text{PDR} = 100 \times \frac{\displaystyle\sum_{i=1}^{N_{\text{STA}}} rx_i}
                                {N_{\text{STA}} \times P}
  \label{eq:pdr}
\end{equation}

Onde $N_{\text{STA}} = 5$ é o número de STAs por experimento e $P = 100$ é o número de pacotes enviados por cada fluxo de medição. Esse denominador é idêntico para AODV e OLSR, tornando o PDR diretamente comparável entre protocolos.

Complementar ao PDR, a taxa de perda de pacotes é definida como:
\begin{equation}
  \text{Loss} = 100 \times \frac{N_{\text{STA}} \times P - \displaystyle\sum_{i} rx_i}
                                 {N_{\text{STA}} \times P}
  \label{eq:loss}
\end{equation}

\paragraph{Latência fim-a-fim (E2E).}
%Para cada fluxo $i$, o FlowMonitor acumula a soma dos atrasos individuais (\texttt{delaySum\textsubscript{i}}) em nanossegundos.
Além disso, mediu-se a latência média ponderada pelo número de pacotes recebidos por fluxo, para refletir o impacto real sobre os dados efetivamente entregues:
\begin{equation}
  \bar{d} = \frac{\displaystyle\sum_{i} \frac{\texttt{delaySum}_i}{rx_i} \cdot rx_i}
                 {\displaystyle\sum_{i} rx_i}
           = \frac{\displaystyle\sum_{i} \texttt{delaySum}_i}
                  {\displaystyle\sum_{i} rx_i}
  \label{eq:delay}
\end{equation}
O resultado é convertido de nanosegundos para milissegundos. Esta formulação equivale à métrica de \textit{time delay} de \cite{turlykozhayeva2024evaluating}: $\text{delay}_i = |EndTime_i - StartTime_i|$.

\paragraph{Jitter.}
O NS-3 calcula o \texttt{jitterSum} de acordo com a definição RFC~3550:
\begin{equation}
  Jitter = \bigl|(D_{i+1} - D_i) - (S_{i+1} - S_i)\bigr|
  \label{eq:jitter}
\end{equation}
Onde $S_i$ e $D_i$ são os instantes de envio e recepção do pacote $i$.A média ponderada por fluxo é calculada de forma análoga à latência (Equação~\ref{eq:delay}).

\paragraph{Throughput.}
O throughput agrega os bytes UDP efetivamente recebidos em todos os fluxos
de medição (wiredSta$\to$backhaul), dividido pela janela de medição~$W$:
\begin{equation}
  \text{Throughput} = \frac{8 \times \displaystyle\sum_{i} \texttt{rxBytes}_i}
                           {W} \quad [\text{kbps}]
  \label{eq:throughput}
\end{equation}
onde $W = T_{\text{duration}} - T_{\text{warmup}}$ é a duração efetiva da
janela de medição (Tabela~\ref{tab:topologias}).
O denominador exclui o período de aquecimento, que seria inflado para OLSR
(que não gera tráfego durante a convergência inicial) e deflado para AODV
(que gera RREQ), garantindo comparação na mesma base temporal.

\paragraph{Hops médios.}
O número médio de saltos por pacote é calculado como:
\begin{equation}
  \bar{h} = \frac{\displaystyle\sum_{i} \texttt{timesForwarded}_i}
                 {\displaystyle\sum_{i} rx_i}
  \label{eq:hops}
\end{equation}
Este indicador quantifica o overhead de roteamento: valores elevados indicam
caminhos subótimos ou retransmissões excessivas causadas pelo protocolo.

Cada par de simulações (AODV, OLSR) utiliza sementes aleatórias distintas e documentadas, seguindo esquema determinístico: para o ponto de escala $k$, a semente AODV é fixada em $100k$ e a semente OLSR em $100k+1$. Por exemplo, no ponto aleatório de 50 nós, emprega-se AODV-seed200 e OLSR-seed201. Este esquema garante que as posições dos nós nas topologias aleatórias sejam geradas deterministicamente a partir dessas sementes, assegurando que a topologia física seja idêntica entre protocolos para um mesmo ponto de escala.

Todos os scripts de geração de configuração, execução de simulação e análise de resultados encontram-se disponíveis no repositório de código aberto \cite{caburlao_github}, permitindo que qualquer pesquisador reproduza exatamente os experimentos aqui reportados.

\section{Resultados e Discussão}



\section{Conclusão e Próximos Passos}



\bibliographystyle{sbc}
\bibliography{Mesh}

\end{document}
